\documentclass[11pt,a4paper]{article}
\usepackage[margin=2.2cm]{geometry}
\usepackage{booktabs,tabularx,array,enumitem,xcolor,hyperref,fancyhdr}
\hypersetup{colorlinks=true,urlcolor=blue,linkcolor=black}
\definecolor{osdagGreen}{HTML}{91B014}
\newcommand{\file}[1]{\texttt{\detokenize{#1}}}
\setlist{nosep,leftmargin=1.4em}
\pagestyle{fancy}\fancyhf{}
\fancyhead[L]{OsdagBridge report generator -- technical changes}
\fancyfoot[C]{\thepage}
\renewcommand{\headrule}{\color{osdagGreen}\hrule height 0.8pt}
\setlength{\parskip}{4pt}\setlength{\parindent}{0pt}
\newcolumntype{Y}{>{\raggedright\arraybackslash}X}

\begin{document}
{\LARGE\bfseries OsdagBridge LaTeX Report Generator\\[2pt] Refactoring \& Enhancements}\\[4pt]
{\large Technical changes document}\par
\vspace{4pt}\textcolor{osdagGreen}{\rule{\textwidth}{1pt}}

\section*{1. Summary}
The PDF design report was reviewed, refactored and extended in
\file{src/osdagbridge/core/reports/}. All six requested items are implemented. In addition, the generated PDF
was \emph{read page by page as a reviewer would} and every defect found was fixed (\S3). Results were verified
by generating the report from real design runs (same input data) with the original code (\file{dev}, commit
\texttt{d74f56d}) and with the changed code, and by an automated QA matrix over five bridge geometries.

\begin{center}
\begin{tabularx}{\textwidth}{@{}Y r r@{}}
\toprule
\textbf{Measure (sample design: 30\,m span, 7.5\,m carriageway, 4 girders)} & \textbf{Before} & \textbf{After}\\
\midrule
LaTeX errors in the compile log & 139 & \textbf{0}\\
Long tables whose header repeats on every page & 1 of 49 & \textbf{51 of 51}\\
Body text inside the footer band (automated check) & 1 (``G4'') & \textbf{0}\\
Executive summary contradicting the results & yes (see \S3) & \textbf{no}\\
Over-wide tables (Overfull \verb|\hbox| warnings) & 60 & 47\\
Graphics embedded (charts, drawings, plots, logo) & 4 & 13\\
Second design (42\,m span, 12\,m carriageway): LaTeX errors / footer overlaps & 152 / 1 & \textbf{0 / 0}\\
QA matrix, 5 designs (20--50\,m span): errors / layout problems & -- & \textbf{0 / 0}\\
Unit tests in \file{tests/unit} & 3 (placeholders) & 78\\
\bottomrule
\end{tabularx}
\end{center}

\textbf{Note on the technology.} The task text mentions PyLaTeX. In this code base the report is assembled
from raw LaTeX strings by the chapter generators (\file{chap1.py}--\file{chap9.py}) and compiled with
\texttt{pdflatex}. The refactor keeps that architecture; the equivalent LaTeX constructs
(\verb|longtable|, \verb|\endhead|, \verb|figure|) are produced by helper functions in the new style module.

\section*{2. The six requested items}
\begin{tabularx}{\textwidth}{@{}>{\raggedright\arraybackslash}p{3.4cm} Y@{}}
\toprule
\textbf{Requirement} & \textbf{What was done}\\
\midrule
1. Repeated table headers &
\verb|styles.longtable_head()| emits the caption, header row, \verb|\endfirsthead|, a ``(continued)'' caption
and header row, and \verb|\endhead|. 49 hand-written tables were converted; the one table with an existing
two-row header (\file{chap4.py}) already repeated correctly. Before, one table repeated its header and a
second (\file{chap2.py}) had an \emph{empty} \verb|\endfirsthead\endhead| pair.\\[3pt]
2. Footer overlap and spacing &
Root causes in \S4. Fixed in \file{styles.py}, \file{chap5.py}. The new checker
\file{tools/report\_layout\_check.py} finds the defect automatically (the original report fails it, the new
one passes on every design tested).\\[3pt]
3. Live-load tables (Ch.\ 3) &
Three tables: \emph{Vehicular (IRC 6)} (one row per vehicle class selected in the UI, with its IRC~6
reference and impact factor), \emph{Associated Vehicular Parameters} (design lanes, braking force,
centrifugal force, eccentricity, with a \textbf{Unit} column) and \emph{Footpath}.\\[3pt]
4. Utilization-ratio charts (Sec.\ 5.5) &
\file{charts.py}: governing UR per primary element and UR of every check; green for UR $\leq 1.0$, red
above, red dashed line at UR\,=\,1.0. Table 5.22 is also heat-shaded (green $\leq0.90$, amber $\leq1.0$,
red $>1.0$).\\[3pt]
5. Quantity charts (Ch.\ 7) &
Steel tonnage (girders / cross bracing / end diaphragms) and concrete volume vs reinforcement weight
(twin axes), from the BOQ values.\\[3pt]
6. Central style module &
\file{styles.py} (\S5). A unit test fails if a chapter file hard-codes colours, paddings, table setup or
\texttt{longtable} boilerplate.\\
\bottomrule
\end{tabularx}

\section*{3. Defects found by reading the report (all fixed)}
\begin{enumerate}
\item \textbf{The executive summary contradicted the results.} \emph{Overall Design Status: Fail} was printed
next to the hard-coded lines ``Girder design pass, Cross bracing design pass, \dots'' (\file{executive_summary.py}).
It now shows a verdict banner, a per-element outcome table and a chart, all derived from the same UR records
as Table 5.22 (\file{verdict.py}). An element with no data is ``Not evaluated'', never ``Pass''.
\item \textbf{Wrong units for girder dimensions.} \texttt{design()} converts section dimensions to metres in
the input dictionary, and the report printed them with an ``mm'' suffix (depth \emph{1.67 mm} instead of
1670\,mm). This also affects reports produced from the GUI. Fixed by \verb|_render_mm()|.
\item \textbf{Table 4.2 (Reactions at Supports) was always empty:} its rows were literally empty strings in the
template. It is now filled from the analysis summary (key load cases first).
\item \textbf{Non-deterministic table order.} Table 4.1 listed load cases in Python \texttt{set} order, which
changes between runs; the same design produced differently ordered reports. A stable, engineering-sensible
order is now used (\verb|load_case_sort_key|).
\item \textbf{Missing values:} number of traffic lanes (the report read a UI counter that is not in the design
dictionary; now \verb|IRC6_2017.table_6|), bare ``kN'' cells for seismic forces the software never computes
(now ``---''), missing units for girder spacing and deck thickness in the summary.
\item \textbf{Layout polish:} justified text in narrow table columns (``Re-strained'', ``Spac- ing'') now
ragged-right; placeholder boxes no longer run off the page; table headers capitalised consistently.
\item \textbf{A regression of my own, found and reverted:} top-aligning \texttt{multirow} labels
(\S4) also moved the two-row header of Table 4.1 out of its cell. The change now applies to data rows only.
\end{enumerate}

\section*{4. Root causes of the layout defects}
\begin{enumerate}
\item \textbf{``G4'' in the footer.} Girder labels are \verb|\multirow{n}{*}| cells, vertically centred over
their group of rows. When the page broke \emph{inside} a group, the centred label was drawn below the last
row on the page, i.e.\ over the footer. \texttt{longtable} cannot reliably be told to keep such a group
together (\verb|\\*|, \verb|\penalty10000| and a larger \verb|LTchunksize| were tried, no effect), so the label
is top-aligned (\verb|\multirow[t]|) and always stays in the first row of its group.
\item \textbf{Footer note.} The ``* Software default value'' note was lifted into the body with
\verb|\vspace{-20pt}|. It now sits in the footer below the rule; \texttt{footskip} is 40\,pt.
\item \textbf{\texttt{\textbackslash needspace} before \texttt{longtable}.} It stranded a table head at the
bottom of a page and printed a duplicated ``(continued)'' header on the next. Removed for long tables.
\item \textbf{LaTeX errors (139).} Mostly ``Misplaced \verb|\noalign|'': a \verb|\caption| followed by
\verb|\hline| with no row terminator. \verb|longtable_head()| always emits \verb|\caption{..}\\|. A
``less-or-equal'' sign (U+2264) was unsupported by \texttt{pdflatex}; \verb|_tex()| now maps common Unicode
symbols to LaTeX.
\item \textbf{Table width.} Column widths exceeded the 15.9\,cm text width. \verb|fit_colspec()| scales the
\texttt{cm} widths to fit, using the cell padding of the table's density.
\item \textbf{Shaded UR cells.} \verb|\cellcolor| (colortbl) paints over a plain \verb+|+ border in the
last column. The summary table therefore ends its column spec with a rule drawn as a column separator
(\verb|RIGHT_RULE|), which is drawn on top of the colour; checked at 300\,dpi (every row has a black border
pixel at the right edge).
\end{enumerate}

\section*{5. Style system and new modules}
\begin{tabularx}{\textwidth}{@{}>{\ttfamily}l Y@{}}
\toprule
\textrm{\textbf{Module}} & \textbf{Responsibility}\\
\midrule
styles.py & single source of truth: DOCUMENT (class, margins, footskip, spacing), COLORS (palette incl.\
status and heat-map colours), TABLES and TABLE\_DENSITY (padding, rules, row height), HEADER\_FOOTER, CHARTS;
builders \verb|preamble()|, \verb|longtable_head()|, \verb|fit_colspec()|, \verb|ur_cell()|,
\verb|verdict_banner()|, \verb|figure_block()|, status helpers\\
verdict.py & per-element and overall pass/fail from one list of UR records (used by summary, charts, tables)\\
charts.py & UR and quantity charts (matplotlib Agg, safe in the report worker thread)\\
drawings.py & to-scale cross-section, plan and elevation schematics from the design inputs; used only when the
GUI did not supply a CAD view (e.g.\ command-line reports), so no empty ``PLACEHOLDER'' boxes\\
\bottomrule
\end{tabularx}

\textbf{Rationale.} One module owns every visual decision, so a global restyle is a one-line edit. Builders are
plain functions returning LaTeX strings: easy to unit-test, no new dependency. Summary, charts and Table 5.22
cannot disagree because they share one data source.

\section*{6. Files}
\begin{tabularx}{\textwidth}{@{}>{\ttfamily}l Y@{}}
\toprule
\textrm{\textbf{File}} & \textbf{Change}\\
\midrule
reports/styles.py, verdict.py, charts.py, drawings.py \textrm{(new)} & see \S5\\
reports/report\_generator.py & preamble from styles; chapter 5 built first so the summary can use its UR records; chart images and fallback drawings\\
reports/chap2.py -- chap5.py, chap7.py, chap8.py & table helpers, units, Table 4.2, ordering, live-load tables, heat map\\
reports/executive\_summary.py & verdict-driven summary (no hard-coded outcomes)\\
reports/report\_utils.py & Unicode mapping, \verb|_render_mm|, lane count, load-case order, placeholder spacing\\
utils/codes/irc6\_2017.py & braking-force fix (\S7)\\
tools/report\_layout\_check.py \textrm{(new)} & footer / overlap checker for generated PDFs\\
tools/generate\_sample\_report.py, report\_qa\_matrix.py \textrm{(new)} & headless design + report in one command; QA over several geometries\\
tests/unit/test\_report\_*.py \textrm{(new)} & 75 new unit tests\\
\bottomrule
\end{tabularx}

\section*{7. Defects found in existing calculation code}
\begin{itemize}
\item \textbf{Fixed -- \texttt{IRC6\_2017.cl\_211\_2\_braking\_force()}} raised \texttt{UnboundLocalError} for
\emph{every} lane count, so the report always showed ``N/A''. It also works with wheel loads in newtons while
its docstring and the report assumed tonnes (a naive fix would print 1\,066\,296\,kN). Initialisation and unit
are corrected; one Class~A train gives 108.69\,kN (20\,\% of 55.4\,t), checked by hand.
\item \textbf{Not changed -- methodology.} The braking-force method (20\,\% of one Class~A train for one or two
lanes) is the software's existing logic and should be reviewed by a structural engineer. Centrifugal force is
not computed (no curvature input); the report states ``Not considered (straight alignment)''.
\end{itemize}

\section*{8. Verification}
\textbf{Method.} The Osdag design pipeline is run headlessly (Qt offscreen, the code path of the \emph{Design}
button; sample project: Mumbai, with weather data) and the report is generated from the resulting data, with
the original and the changed code. The analysis plots are attached exactly as the desktop app does
(\verb|capture_report_figures|).

\begin{center}
\begin{tabular}{rrrrrr}
\toprule
Span (m) & Carriageway (m) & Pages & LaTeX errors & Layout problems & Time (s)\\
\midrule
20 & 7.5 & 56 & 0 & 0 & 31\\
30 & 7.5 & 56 & 0 & 0 & 36\\
36 & 9.0 & 56 & 0 & 0 & 40\\
42 & 12.0 & 63 & 0 & 0 & 77\\
50 & 15.0 & 66 & 0 & 0 & 124\\
\bottomrule
\end{tabular}
\end{center}

\textbf{Limits.} Chapter 6 3D CAD views and the 3D view of the executive summary need the GUI viewer, so they
remain placeholders in command-line reports; the new schematics replace only the cross-section, plan and
elevation. In the ``Before'' PDFs those slots are empty because the original code has no fallback (the
desktop app would supply CAD images). End-diaphragm quantities are not produced by the BOQ module (``N.A.''
in the chart). The weather inputs are sample values. The change was verified with a standard TeX Live
installation, not with the Windows \texttt{osdag\_latex\_env} bundle; it adds no LaTeX package.

\section*{9. Reproduce}
\begin{verbatim}
python -m pytest tests/unit                          # 78 tests
python tools/generate_sample_report.py --span 30 --carriageway 7.5 --out out
python tools/report_qa_matrix.py --out qa_out        # 5 designs, exits 1 on any problem
python tools/report_layout_check.py out/Report.pdf   # needs poppler pdftotext
\end{verbatim}
\end{document}
